\documentclass[10pt,a4paper]{amsart}
\usepackage[margin=19mm]{geometry}
\usepackage{amsmath,amssymb,mathtools}
\usepackage[scr=rsfs,cal=euler]{mathalfa}
\usepackage{xfrac}
\usepackage[unicode,hidelinks,bookmarks=false]{hyperref}
\newcommand{\ie}{i.e.\ }
\newcommand{\cf}{cf.\ }
\newcommand{\resp}{respectively\ }
\newcommand{\Subsection}{\subsection}
\providecommand{\nobreakdash}{\nobreak\hbox{-}}
\setlength{\emergencystretch}{2em}
\multlinegap0pt
\usepackage{bm}
\usepackage{bbm}
\usepackage{dsfont}
\usepackage{xspace}
\usepackage{cancel}
\usepackage{mathtools}
\usepackage{amsthm}
\makeatletter
\@ifundefined{thm}{\newtheorem{thm}{Theorem}}{}
\makeatother
\title[Inverse Scattering for Dirac Equations Arising in Waveguide ]{Inverse Scattering for Dirac Equations Arising in Waveguide Arrays}
\author{John C. Schotland, Shenwen Yu}
\date{Source: arXiv:2605.01230v1 (CC BY 4.0)}
\begin{document}
\maketitle

\noindent\emph{Source and licence.} Inverse Scattering for Dirac Equations Arising in Waveguide Arrays. Version arXiv:2605.01230v1 (CC BY 4.0), distributed under the Creative Commons Attribution 4.0 International licence, as recorded in the arXiv licence metadata for this exact version and shown on the abstract page https://arxiv.org/abs/2605.01230v1. Full rights notice: licenses/arxiv-2605.01230-CC-BY-4.0.txt.\par\smallskip

\noindent\textit{Editorial scope.} Counted unit: the Thm whose source label is \texttt{iter of chiral} in the article (source0.tex lines 156--167, source label \texttt{iter of chiral}), delivered together with the article's own complete original proof of it (source0.tex lines 169--197, one \texttt{proof} environment of 29 lines). No mathematics is written, repaired or re-derived: statement and proof are the article's own text line for line. Required context retained uncounted: chiral model (lines 103--114). Every definition, display and label of those lines is retained unchanged. Omitted from this package and not redistributed: 1--102, 115--155, 168--168, 198--1479. Nothing inside a retained excerpt is omitted or altered: every retained body is delivered exactly as the article's lines are written, so no omission receipt of any kind accompanies this package. Citations: the retained excerpts carry no citation command at all, so no bibliography entry of the article is transcribed and no reference is redistributed. Reference adaptation: none was needed and none was made; every reference inside the delivered unit resolves inside it, so no unresolved reference or citation is produced. Printed slips kept literal: no printed word, symbol, hypothesis or number of the counted statement or of its proof was altered. Layout and class: the article's own class is \texttt{amsart} with packages including fullpage, amsmath, amssymb, amsfonts, amsthm, bm, bbm, comment, caption, placeins, graphicx, subcaption, booktabs, float; the wrapper is the lane's pinned \texttt{amsart} editorial wrapper at 10pt on A4, which restates the theorem-like environments the retained text needs and adds the article's own macro definitions that the retained text needs (none), with extra wrapper packages (bm, bbm, dsfont, xspace, cancel, mathtools). The delivered numbering is the unit's own. Assets: no retained span contains a figure environment, a graphics-inclusion command, a PDF-page inclusion, a table or a TikZ picture; no image file is copied into this package, the package declares no asset dependency and no image file is redistributed with it. Licence: version arXiv:2605.01230v1 of this article is distributed under the Creative Commons Attribution 4.0 International licence, as recorded in the arXiv licence metadata for this exact version and shown on the abstract page https://arxiv.org/abs/2605.01230v1. A full rights notice is delivered at licenses/arxiv-2605.01230-CC-BY-4.0.txt. Characters: every retained excerpt is pure ASCII, so the delivered engine, XeLaTeX with its default Latin Modern text font, reports no missing character.
\begin{equation}\label{chiral model}
\left\{
\begin{aligned}
  &-\mathrm{i}\,\partial_x \psi + \mathrm{i}\,\alpha\,\partial_y \psi + k\bigl[\beta + V(x,y)\bigr]\psi = 0
  && \text{in } \Omega,\\[0.4em]
  &\psi_2(x,0)=0,\quad \psi_1(x,L_y)=0
  && \text{for } 0\le x\le L_x,\\[0.4em]
  &\psi(0,y)=g(y)
  && \text{for } 0\le y\le L_y,
\end{aligned}
\right.
\end{equation}

\begin{thm}
Assume that $V\in L^{\infty}(\Omega)$. Then the initial--boundary value problem \eqref{chiral model} admits a unique solution
\[
\psi\in C\bigl([0,L_x];L^2(0,L_y;\mathbb{C}^2)\bigr)
\]
satisfying the integral equation
\begin{equation}\label{iter of chiral}
    \psi(x)
    = T(x)g
    + \int_0^x T(x-x')B(x')\psi(x')\,dx'.
\end{equation}
\end{thm}

\begin{proof}
Set
\[
x_0=\frac{1}{2k\|V\|_{L^{\infty}(\Omega)}},
\quad
X_0:=C\bigl([0,x_0];L^2(0,L_y;\mathbb{C}^2)\bigr),
\]
equipped with the supremum norm. Define the mapping $\Phi_c:X_0\to X_0$ by
\[
(\Phi_c\psi)(x)
:=
T(x)g
+
\int_0^x T(x-x')B(x')\psi(x')\,dx'.
\]
Using $\|T(x)\|_{L^2\to L^2}=1$ and $\|B(s)\|_{L^2\to L^2}\le k\|V\|_{L^\infty(\Omega)}$, we obtain
\begin{align*}
    \|\Phi_c\psi-\Phi_c\varphi\|_{X_0}
&\leq
\left(\int_0^{x_0}k\|V(s,\cdot)\|_{L^\infty([0,L_y])}\,ds\right)\|\psi-\varphi\|_{X_0}\\
&\leq
 x_0 k\|V\|_{L^\infty(\Omega)}\,\|\psi-\varphi\|_{X_0}.
\end{align*}
By the choice of $x_0$,  we have $ x_0k\|V\|_{L^\infty(\Omega)}=\tfrac12$, hence $\Phi_c$ is a contraction on $X_0$. Existence and uniqueness of the solution on $[0,x_0]$ follow from Banach's fixed point theorem. We then repeat the same argument on the successive subintervals
\[
[x_0,2x_0],\ [2x_0,3x_0],\ \ldots,
\]
thereby obtaining a unique solution on each subinterval and, consequently, a unique solution on the entire interval $[0,L_x]$. 
\end{proof}

\end{document}
